%!TEX program = xelatex
\documentclass[a4paper,12pt]{article}
\usepackage{fontspec}\defaultfontfeatures{Ligatures=TeX}
\usepackage[a4paper,vmargin={3cm,3.5cm},hmargin={3cm,3cm}]{geometry}
%-----------------------------------------------------------------------------%
\usepackage{settings}
%-----------------------------------------------------------------------------%
%%% Title %%%
\title{G6411 Recitation Notes: Distributions of Hypothesis Testing}
\author{Jesse Chieh Chen\\\href{mailto:cc5229@columbia.edu}{cc5229@columbia.edu}}
\date{\today}
%-----------------------------------------------------------------------------%

\begin{document}

\appendix
\section{Distributions of Hypothesis Testing}

The normal, chi-squared, student $t$, and $F$ distributions are closely related.
We define them through standard normal random variables and use these relationships
to derive their limits and an application to regression residuals.
Throughout, degrees of freedom are positive integers.

\subsection{The Normal Distribution}

\begin{definition}[Standard Normal Distribution]
	A random variable $Z$ has the standard normal distribution, written $Z\sim\normaldistribution(0,1)$,
	if its density is
	\begin{align}
		\phi(z)=\frac{1}{\sqrt{2\pi}}\exp\left(-\frac{z^2}{2}\right),
		\qquad z\in\reals.
	\end{align}
	It has mean zero and variance one.
	More generally, $Y=\mu+\sigma Z\sim\normaldistribution(\mu,\sigma^2)$ for $\sigma>0$,
	so $(Y-\mu)/\sigma\sim\normaldistribution(0,1)$.
\end{definition}

\begin{remark}[Normal Vectors]
	If $Z=(Z_1,\ldots,Z_n)\T\sim\normaldistribution(0,I_n)$,
	then its components are independent standard normal random variables.
	For any fixed matrix $A$, $AZ\sim\normaldistribution(0,AA\T)$.
	In particular, if $U$ is orthogonal, then
	$U\T Z\sim\normaldistribution(0,I_n)$:
	rotating a standard normal vector preserves its distribution.
\end{remark}

\subsection{The Chi-Squared Distribution}

\begin{definition}[Chi-Squared Distribution]
	Let $Z_1,\ldots,Z_r$ be independent standard normal random variables.
	Then
	\begin{align}
		V=\sum_{j=1}^rZ_j^2\sim\chi_r^2.
	\end{align}
	We say that $V$ has a chi-squared distribution with $r$ \emph{degrees of freedom}.
	Here the degrees of freedom count the independent standard normal components being squared and added.
\end{definition}

\begin{proposition}\label{prop-chi-square-concentration}
	If $T_r$ follows $\chi_r^2$, then
	\begin{align}
		\E(T_r)=r,\qquad \var(T_r)=2r,
		\qquad \frac{T_r}{r}\pto 1\quad\text{as }r\to\infty.
	\end{align}
\end{proposition}
\begin{proof}
	For $Z\sim\normaldistribution(0,1)$, $\E(Z^2)=1$ and $\E(Z^4)=3$,
	so $\var(Z^2)=2$.
	Independence gives the stated mean and variance of the sum.
	Thus $\E(T_r/r)=1$ and $\var(T_r/r)=2/r\to0$;
	Chebyshev's inequality gives the convergence in probability.
\end{proof}

\subsection{The Student \texorpdfstring{$t$}{t} Distribution}

\begin{definition}[Student $t$ Distribution]
	Let $Z\sim\normaldistribution(0,1)$ and $V\sim\chi_\nu^2$ be independent.
	Then
	\begin{align}
		T=\frac{Z}{\sqrt{V/\nu}}\sim t_\nu.
	\end{align}
	We say that $T$ has a $t$ distribution with $\nu$ degrees of freedom.
	Its degrees of freedom come from the chi-squared variable in the denominator.
\end{definition}

\begin{remark}
	The $t$ distribution is symmetric about zero, like the standard normal,
	but has heavier tails because its denominator is random and can be small.
	The independence of $Z$ and $V$ is essential for the stated distribution.
\end{remark}

\begin{remark}
	A $t$ distribution with one degree of freedom is a Cauchy distribution.
\end{remark}

\begin{proposition}[Normal Limit of the $t$ Distribution]
	If $T_\nu$ follows $t_\nu$, then $T_\nu\dto\normaldistribution(0,1)$ as $\nu\to\infty$.
\end{proposition}
\begin{proof}
	Write $T_\nu=Z/\sqrt{V_\nu/\nu}$ with $Z\sim\normaldistribution(0,1)$
	independent of $V_\nu\sim\chi_\nu^2$.
	By \autoref{prop-chi-square-concentration} and the continuous mapping theorem,
	$\sqrt{V_\nu/\nu}\pto1$.
	Slutsky's theorem therefore gives $T_\nu\dto\normaldistribution(0,1)$.
\end{proof}

\subsection{The \texorpdfstring{$F$}{F} Distribution}

\begin{definition}[$F$ Distribution]
	Let $U\sim\chi_r^2$ and $V\sim\chi_\nu^2$ be independent.
	Then
	\begin{align}
		F=\frac{U/r}{V/\nu}\sim F_{r,\nu}.
	\end{align}
	The first degree-of-freedom parameter, $r$, belongs to the numerator;
	the second, $\nu$, belongs to the denominator.
	Each chi-squared variable is divided by its own degrees of freedom.
\end{definition}

\begin{remark}[A Squared $t$ Distribution follows an $F$ Distribution]
	If $T\sim t_\nu$, then $T^2\sim F_{1,\nu}$
	since $T^2={Z^2}/(V/\nu)$ where $Z^2\sim\chi_1^2$ is independent of $V$.
\end{remark}

\begin{proposition}[Chi-Squared Limit of the $F$ Distribution]
	For any fixed $r$, if $T_\nu$ follows $F_{r,\nu}$, then $rT_\nu\dto\chi_r^2$ as $\nu\to\infty$.
\end{proposition}
\begin{proof}
	Write $T_\nu=(U/r)/(V_\nu/\nu)$ with independent
	$U\sim\chi_r^2$ and $V_\nu\sim\chi_\nu^2$.
	Since $V_\nu/\nu\pto1$, Slutsky's theorem gives the result.
	The numerator degrees of freedom stay fixed in this limit.
\end{proof}

\subsection{Normal Quadratic Forms}

\begin{proposition}\label{thm-normal-quadratic-form}
	Let $Z\sim\normaldistribution(0,I_n)$ and let $M$ be a fixed $n\times n$ symmetric idempotent matrix
	with rank $r\ge1$. Then
	\begin{align}
		Z\T MZ\sim\chi_r^2.
	\end{align}
\end{proposition}
\begin{proof}
	Since $M$ is symmetric, it has an orthogonal eigendecomposition $M=UDU\T$.
	Idempotence implies that each eigenvalue satisfies $\lambda^2=\lambda$,
	so each is either zero or one.
	Since the rank is $r$, we can order the eigenvalues so that
	\begin{align*}
		D=\begin{pmatrix}I_r&0\\0&0\end{pmatrix}.
	\end{align*}
	Let $W=U\T Z\sim\normaldistribution(0,I_n)$. Then
	\begin{align*}
		Z\T MZ&=W\T DW=\sum_{j=1}^rW_j^2\sim\chi_r^2.
		\qedhere
	\end{align*}
\end{proof}

\begin{remark}
	A symmetric idempotent matrix is an orthogonal projection.
	Its rank counts the normal components retained by the projection,
	which explains the degrees of freedom in the theorem.
	If the rank is zero, the quadratic form is identically zero.
\end{remark}

\begin{example}[Regression Residuals]
	Consider $Y=X\beta+e$, where $X$ is $n\times k$ with full column rank, $n>k$, and
	\begin{align*}
		e\given X\sim\normaldistribution(0,\sigma^2I_n),\qquad \sigma^2>0.
	\end{align*}
	The \gls{ols} residual vector is
	\begin{align*}
		\widehat e=MY=Me,\qquad M=I_n-X(X\T X)\inv X\T.
	\end{align*}
	Conditional on $X$, the matrix $M$ is fixed, symmetric, and idempotent,
	with rank $n-k$. Also $e/\sigma\given X\sim\normaldistribution(0,I_n)$.
	Since $M\T M=M$, \autoref{thm-normal-quadratic-form} gives
	\begin{align}
		\frac{\widehat e\T\widehat e}{\sigma^2}
		=\left(\frac{e}{\sigma}\right)\T M\left(\frac{e}{\sigma}\right)
		\given[\Big] X\sim\chi_{n-k}^2.
	\end{align}
	The conditional distribution does not depend on $X$, so the same distribution holds unconditionally.
	Estimating $k$ coefficients removes $k$ dimensions, leaving $n-k$ residual degrees of freedom.
	Here $k$ includes the intercept if one is estimated.
\end{example}


\end{document}
